% !PRINTED 618
\weberpage{618}
%% [page furniture, not body text -- source-confirmed 2026-08-10]
%% 618 \quad Twenty-Third Chapter. \quad \S\,168.

The first of the conditions (9) requires that $p$ be properly representable by the principal form of the discriminant $D$ (\S\,122), the second that
\weq{10}{p \equiv 1 \pmod{m}}
hold.

In order that a prime ideal $\mathfrak{p}$ of the first degree in $\Omega$, which does not divide $S$, be contained in $\bar{E}A$, it is necessary and sufficient that

1. $\mathfrak{p}$ be a principal ideal, i.\,e. an existing number $\pi$, and that this number $\pi$ be contained in $A$, i.\,e. that
\[
\text{2. } N(\pi) = p \equiv 1 \pmod{m}
\]
hold. These conditions agree exactly with the conditions (9), and consequently $\mathfrak{K}(D, \varrho)$ is \gesperrt{the class field of the group $A$}.

The degree of the field $\mathfrak{K}(A)$ is therefore not merely at most, but exactly equal to $h$.

The number $(Z, R_m)$ is $=1$ if none of the characteristic primes of \S\,108 divides $m$, and then the degree of the field $\mathfrak{K}(A)$ is equal to $H'\varphi(m)$. If, conversely, we include in $m$ all the characteristic primes and prime powers, then $R_m = R$ becomes \gesperrt{the group of absolute norm residues}, and we have
\[
(Z, R) = 2g,
\]
if $g$ is the number of genera of the discriminant $D$ (\S\,113, 7.). If, then, $H'_g$ denotes the number of classes of one genus, then $H' = g H'_g$, and formula (7) gives:
\weq{11}{h = \tfrac{1}{2} H'_g \varphi(m).}

By \S\,138 the class function $H_{-D}(u)$ splits, through the adjunction of square roots, into as many factors as there are genera of classes of forms, and the degree of each of these factors is $H'_g$.

Now if the class equation were still further decomposable, through the adjunction of any roots of unity whatever, than according to the genera, then one could assume $m$ in such a way that these roots of unity too, and also $\sqrt{\Delta}$, were contained in the field $\mathfrak{R}(\varrho)$.

If then $h''$ is the lowest degree to which the class equation can be reduced by the adjunction of roots of unity, then the absolute degree of the field $\mathfrak{K}(D, \varrho)$ is at most equal to $h''\varphi(m)$
%CONT>>
